% 由 scripts/make_paper_appendix.py 自动生成，请勿手改。

\section{Prompt templates}
\label{app:prompts}
All prompts are rendered from a single template table in \texttt{src/prompts.py};
the MCQ templates are selected automatically when the question contains at least two
options of the form ``A. \dots''. Every template demands the final answer on its own
line prefixed with \texttt{ANSWER:}, which is what the extractor looks for first.

\paragraph{Cot / free-response}
{\footnotesize\begin{verbatim}  
Solve the following physics problem step by step. Show all work.
At the end, give the final numeric answer on its own line, prefixed with "ANSWER:".
Problem: <problem statement>
\end{verbatim}}
\paragraph{Cot / multiple-choice}
{\footnotesize\begin{verbatim}  
Solve the following multiple-choice physics problem step by step. Show all work.
At the end, give the chosen option letter on its own line, prefixed with "ANSWER:".
Problem: <problem statement>
\end{verbatim}}
\paragraph{Direct / free-response}
{\footnotesize\begin{verbatim}  
Answer the following physics problem directly with a single number (no explanation).
Give the final answer on its own line, prefixed with "ANSWER:".
Problem: <problem statement>
\end{verbatim}}
\paragraph{Direct / multiple-choice}
{\footnotesize\begin{verbatim}  
Answer the following multiple-choice physics problem.
Give the chosen option letter on its own line, prefixed with "ANSWER:".
Problem: <problem statement>
\end{verbatim}}
\paragraph{Structured / free-response}
{\footnotesize\begin{verbatim}  
Solve the following physics problem step by step.
Before solving, first write down the knowns and unknowns explicitly, then derive the solution.
Show all work. At the end, give the final numeric answer on its own line, prefixed with "ANSWER:".
Problem: <problem statement>
\end{verbatim}}
\paragraph{Structured / multiple-choice}
{\footnotesize\begin{verbatim}  
Solve the following multiple-choice physics problem step by step.
Before solving, first write down the knowns and unknowns explicitly, then compare the options.
At the end, give the chosen option letter on its own line, prefixed with "ANSWER:".
Problem: <problem statement>
\end{verbatim}}
\paragraph{Chat template}
The rendered prompt is then wrapped by the tokenizer's chat template
(\texttt{apply\_chat\_template}), so the model sees its own instruction format.
Whether a chat template was applied is recorded in every result file as
\texttt{meta.chat\_template}.
Example problem used in this paper: \texttt{A 2.0 kg block slides down a 30 degree incline at 3.0 m/s. ...}

\section{Result-file schema}
\label{app:schema}
Every evaluation writes one JSON file. The top level records the run identity and
provenance; \texttt{detail} holds one record per problem, which is what the audit
tables are built from. Minimal schema of \texttt{results/eval\_*.json}:

\begin{verbatim}
  "model": string,
  "benchmark": string,
  "mode": string,
  "n": int,
  "n_scorable": int,
  "n_unscored": int,
  "n_unparsed_pred": int,
  "overall_acc": string,
  "per_domain_acc": object,
  "detail": [ ... 323 records ... ],
  "wall_time_s": float,
  "max_new_tokens": int,
  "data_composition": object,
  "prompt_fingerprint": string,
  "meta": object,
  # detail[i]:
  #   "id": string,
  #   "question": string,
  #   "domain": string,
  #   "subset": string,
  #   "ref": string,
  #   "pred_raw": string,
  #   "pred_val": float,
  #   "correct": bool,
  #   "status": string,
  #   "scorable": bool,
  #   "n_pred_chars": int,
\end{verbatim}

Two fields deserve emphasis. \texttt{scorable} marks whether the \emph{reference}
answer could be parsed; rows with \texttt{scorable=false} are excluded from every
denominator and reported as \texttt{n\_unscored}. \texttt{status} is one of
\texttt{correct}, \texttt{incorrect}, \texttt{unparsed\_pred} or
\texttt{unscorable\_ref}, so a reader can recompute any accuracy in this paper from
the raw file without re-running a model. Sensitivity files
(\texttt{sensitivity\_*\_*.json}) additionally carry \texttt{items[]}, one entry per
perturbed problem with the perturbed question text and both answers, enabling
per-item audit of the perturbation.

\section{Statistical details}
\label{app:stats}
\paragraph{Bootstrap.}
Confidence intervals resample problems with replacement $B$~=~10{,}000 times using
\texttt{random.Random} seeded with 42, and take the 2.5/97.5 percentiles of the
resampled statistic (order statistics of the resampled accuracy, i.e.\ a percentile
interval), the standard non-parametric approach of \citet{efron1993bootstrap}.
The same seed is used for every model, so intervals are exactly
reproducible across machines and do not depend on the NumPy version.

\paragraph{Layer-wise interval and paired test.}
The recovery curve has $n=14$ problems. Each layer's mean gets a
percentile bootstrap interval from the same $B$; to test whether the point estimate's
best layer is meaningful, we bootstrap the \emph{difference} between that layer and
every other layer (2{,}000 resamples) and ask whether the interval excludes zero.
With this test the best layer (layer 8) is \emph{not} separable
from 6 other layers, which is why the paper
claims a band rather than a peak.

\paragraph{Small-sample proportions.}
Proportions computed on few observations are reported with Wilson score intervals
(95\%), which behave better than the normal approximation near 0 and 1. For example the
layer-parse recovery on the correct-baseline group is a proportion of a~\nbCleanCorrect{}-problem sample.

\paragraph{Paired tests.}
Correctness comparisons between two conditions on the same problems use McNemar's test
\citep{mcnemar1947} on the discordant pairs. When the number of discordant pairs is
small we report the exact binomial two-sided $p$-value rather than the asymptotic
$\chi^2$ version.

\paragraph{Annotator agreement.}
We report both unweighted Cohen's $\kappa$ and a linearly weighted $\kappa_w$
\citep{landis1977kappa} that gives partial credit for adjacent categories, plus
per-class agreement. The distinction matters because the categories are heavily
imbalanced (E2 dominates), which depresses $\kappa$ even when most individual
decisions agree.

\begin{table}[h]
  \centering
  \caption{Per-class agreement between the same-family and cross-family judge
           ($n=\nbAnnotated$ samples; counts are per class).}
  \label{tab:perclass}
  \small
  \begin{tabular}{@{}l c c c c@{}}
    \toprule
    \textbf{Class} & \textbf{cross-family count} & \textbf{same-family count} &
      \textbf{both} & \textbf{agreement given same-family label} \\
    \midrule
    E1 & 129 & 41 & 20 & 16\% \\
    E2 & 303 & 414 & 267 & 88\% \\
    E3 & 46 & 40 & 11 & 24\% \\
    N & 17 & 0 & 0 & 0\% \\
    \bottomrule
  \end{tabular}
\end{table}
